\clearpage
\fancyhead[L]{\small INSTRUMENTE.}
\SourceLeaf{440}
\begin{center}
{\large INSTRUMENTE.}
\end{center}

\begin{center}
\noindent\rule{\linewidth}{0.4pt}

Göttingische gelehrte Anzeigen. Nr. 187. S. 1865..1866. 20. November 1820.

\noindent\rule{\linewidth}{0.4pt}
\end{center}

Seine Königliche Hoheit, der Herzog von \textsc{Clarence}, hat der Universität ein Zeichen seines gnädigen Gedenkens gegeben, indem er der Bibliothek eine kostbare Sammlung von Seekarten auf 182 Blättern schenkte. Es sind die Karten, die das Hydrographical Office in London herausgibt und mit seinem Stempel versieht. Sie kommen nicht in den Buchhandel und sind ausschließlich für die Königliche Marine bestimmt. Die Sammlung umfasst die europäischen Gewässer ebenso wie den größten Teil der Küsten Afrikas, Amerikas, Ost- und Westindiens. Viele davon kennt der Herzog, selbst Seemann, aus eigener Anschauung. Damit dürften ihr Wert und unser Dank hinreichend zum Ausdruck kommen, denn ein Verzeichnis der einzelnen Karten würde dem Zweck dieser Blätter nicht entsprechen.

Auch Seine Königliche Hoheit, der Herzog von \textsc{Sussex}, hat der Universität mit zwei Geschenken an die Sternwarte ein Zeichen seines huldvollen Gedenkens gegeben. Es handelt sich um eine von \textsc{Hardy} in London gefertigte Tertienuhr von besonderer Bauart und besonders schöner Ausführung und um einen Apparat, den man ein umgekehrtes Pendel nennen könnte und der zeigt, wie sich eine einfallsreiche Idee praktisch nutzen lässt. Das gewöhnliche Pendel ist oben befestigt und \emph{hängt} an der Feder. Dieses ist unten befestigt und \emph{steht} auf der verhältnismäßig kräftigen Feder. Die Pendelstange ist ein 5 Zoll langer Zylinder, auf dessen oberes Ende ein kleines kugelförmiges Gewicht aufgeschraubt ist. Oben läuft die Stange in eine feine Spitze aus, die sich über einem einen Zoll großen, in 20 Teile unterteilten Gradbogen bewegt. Jeder Teil entspricht demnach eigentlich etwas mehr als 34 Minuten. Die Versuche an der Sternwarte haben jedoch ergeben, dass schon eine Neigung des Apparats um 57 Sekunden einen Ausschlag um einen Teil bewirkt. Das Pendel ist also ebenso empfindlich wie ein gewöhnliches von 15 Fuß Länge. Da man das Kügelchen höher oder tiefer schrauben kann, lässt es sich bei Bedarf leicht so abgleichen, dass ein Ausschlag um einen Teil genau einer Minute entspricht. Daraus sieht man, dass der Apparat, mit den nötigen Korrekturschrauben versehen und mit der gebotenen Vorsicht angewandt, in manchen Fällen mit Vorteil an die Stelle einer Libelle oder eines Lots treten kann.

